% !PRINTED 350
\weberpage{350}
%<<CONT

%% [running head, not body text -- commented out 2026-08-07] Page furniture (running
%% head, folio, signature) is excluded from this edition by standing policy.
%% \noindent 350\hfill Eleventh Chapter.\hfill \S\ 95.

With a positive discriminant there exist in every class of forms forms with positive first coefficients, and therefore classes of forms and classes of ideals correspond to one another in a one-to-one manner.

A basis of the principal ideal $\eta$ is $(\eta, \eta\theta)$ and
\[
\left.\begin{aligned}
\theta &= \frac{\sqrt{\Delta}}{2}, & \theta^2 &= \frac{\Delta}{4}, & \Delta &\equiv 0\\
\theta &= \frac{-1+\sqrt{\Delta}}{2}, & \theta^2 &= \frac{\Delta-1}{4} - \theta, & \Delta &\equiv 1
\end{aligned}\right\}\ (\mathrm{mod}\ 4).
\]
If $x, y$ are rational integers and
\[
\eta = x + y\theta,
\]
then this basis is, in the two cases:
\[
\begin{aligned}
(\alpha_1, \alpha_2) &= \left(x + y\theta,\ \frac{\Delta y}{4} + \theta x\right),\\
&= \left(x + y\theta,\ y\,\frac{\Delta-1}{4} + (x-y)\theta\right)
\end{aligned}
\]
and the substitution determinant $A$ [\S\ 87, (1)] is
\[
x^2 - \frac{\Delta}{4}y^2, \quad x^2 - xy + \frac{1-\Delta}{4}y^2 = N(\eta).
\]
The basis $(\alpha_1, \alpha_2)$ is therefore positive if the norm of $\eta$ is positive (otherwise $\alpha_1$ would have to be interchanged with $\alpha_2$). From the basis form
\[
\lambda = \alpha_1 t_1 + \alpha_2 t_2 = \eta(t_1 + \theta t_2)
\]
one obtains
\[
N(\lambda) = N(\eta)\,T,
\]
and thereupon $T$ becomes the \gesperrt{principal form}
\weq{9}{
(1,\,0,\,\tfrac{\Delta}{4}), \quad \left(1,\,-1,\,\frac{+1-\Delta}{4}\right).
}
If $N(\eta)$ were negative, one would have obtained for $T$ the forms
\[
\left(-1,\,0,\,\frac{-\Delta}{4}\right), \quad \left(-1,\,+1\ \frac{-1-\Delta}{4}\right) % [sic] -- no comma printed before the fraction
\]
and these are equivalent to (9) only if the \gesperrt{Pell} equation $t^2 - \Delta u^2 = -4$ is solvable.

\rule{4cm}{0.4pt}
